\chapter{Použité technologie}
\label{sec:technologie}

V této kapitole jsou popsány softwarové nástroje, knihovny a technologie, které byly využity při návrhu a implementaci praktické části této diplomové práce. 
Vzhledem k požadavkům na interaktivitu a vizualizaci formálních modelů byly zvoleny moderní webové technologie, které zajistí vysokou modularitu a přívětivé uživatelské rozhraní.

\section{Vývojové prostředí a sestavení aplikace}

\subsection{Visual Studio Code}
Jako primární vývojové prostředí byl zvolen editor Visual Studio Code (VS Code) od společnosti Microsoft. 
Jedná se o open-source editor, který se stal standardem pro vývoj webových aplikací. 
Jeho hlavní výhodou v kontextu této práce je nativní podpora jazyka TypeScript a Node.js. 
Editor nabízí zabudovanou spolupráci s verzovacím nástrojem Git a integrovaný terminál, přes který je pak možné spouštět lokální vývojový server nebo skripty pro sestavení aplikace, bez nutnosti opouštět okno editoru.

\subsection{GitHub}
Pro verzování zdrojových kódů a zálohování projektu byl využit systém Git v kombinaci s platformou GitHub. 
Využití repozitáře na této platformě zajistilo nejen transparentní historii vývoje pro vedoucího práce, ale také bezpečné a snadno dostupné úložiště pro zdrojové kódy. 

Platforma GitHub nabízí také spoustu rozšíření, příkladem může být služba GitHub Pages, která byla v této práci využita pro kontinuální nasazení (Continuous Deployment) webové aplikace. 
Pro usnadnění přístupu k aplikaci byla k repozitáři připojena vlastní doména. 
Aktuální verze aplikace je tedy dostupná na adrese \texttt{ccs.koberskyj.cz}. 
Využití služby GitHub Pages poskytuje stabilní a ověřené hostingové řešení, které garantuje vysokou dostupnost aplikace pro uživatele.

\subsection{Node.js a npm}
Node.js je běhové prostředí (runtime environment) postavené na enginu V8 z prohlížeče Google Chrome, které umožňuje spouštět JavaScript kód mimo webový prohlížeč. 
V tomto projektu neslouží Node.js pro běh serverové části (backendu), jelikož je aplikace navržena jako čistě klientská, ale poskytuje infrastrukturu pro běh vývojových a sestavovacích nástrojů. 

Společně s Node.js byl využit balíčkovací systém \texttt{npm} (Node Package Manager), který spravuje všechny externí závislosti projektu. 
Pomocí \texttt{npm} byly do projektu integrovány veškeré knihovny zmiňované v této kapitole. 
Správa verzí v souboru \texttt{package.json} navíc zajišťuje deterministická sestavení aplikace na různých počítačích s různým operačním systémem.

\subsection{Vite}
Pro samotné sestavení a běh vývojového serveru byl použit nástroj Vite. 
Oproti starším bundlerům, jako je například Webpack, přináší Vite značné zrychlení vývojového procesu. 
Využívá nativních ES modulů v prohlížeči, což znamená, že při úpravě kódu není nutné překompilovávat celou aplikaci. 

Změny provedené ve zdrojovém kódu se díky funkci Hot Module Replacement (HMR) projevují v prohlížeči téměř okamžitě.
Ve fázi produkčního sestavení pak Vite využívá nástroj Rollup, který kód optimalizuje, minifikuje a připraví pro statické nasazení, například právě na službu GitHub Pages.

\section{Základní webové technologie a uživatelské rozhraní}
Klientská část aplikace musí být, podle požadavků specifikovaných v předchozí kapitole, dynamická a schopná okamžitě reagovat na vstupy uživatele. 
To je klíčové například při interaktivní konstrukci SOS důkazů, kdy je nezbytné okamžitě validovat aplikovatelnost sémantických pravidel. 
Z tohoto důvodu byla zvolena architektura Single Page Application (SPA). 
V tomto modelu se veškeré HTML, JavaScript a CSS načtou pouze jednou a obsah stránky je dynamicky přepisován na základě interakce uživatele bez nutnosti znovunačtení celé webové stránky ze serveru.

\subsection{React.js a TypeScript}
Pro vývoj uživatelského rozhraní byla použita JavaScript knihovna React.js, vyvinutá společností Meta. 
React umožňuje tvorbu komplexních uživatelských rozhraní skládáním malých, znovupoužitelných a izolovaných komponent. 
Deklarativní přístup Reactu výrazně usnadňuje synchronizaci vnitřního stavu aplikace (například informace o aktuálním kroku v SOS důkazu nebo vizuálním stavu LTS grafu) s tím, co uživatel vidí na obrazovce.

Knihovna React byla také doplněna o jazyk TypeScript, který obohacuje standardní JavaScript o statické typování. 

Při návrhu React komponent byly navíc uplatňovány principy metodiky SOLID, což je sada doporučení pro psaní znovupoužitelných komponent. 
Jedná se například o jasné vymezení odpovědnosti komponenty nebo aby byly komponenty otevřené pro vnější rozšíření.
Příkladem takovéto komponenty může být \texttt{CCSViewer}, který umožňuje zvýraznit syntaxi vloženého CCS výrazu a zobrazit jej jako klasický textový element, jakým je například značka \texttt{span}.

\subsection{Tailwind CSS}
Pro stylování aplikace byl využit framework Tailwind CSS. 
Na rozdíl od tradičních CSS frameworků (jako je např. Bootstrap), které poskytují předpřipravené komponenty, je Tailwind \enquote{utility-first} framework. 
Poskytuje nízkoúrovňové třídy (např. pro nastavení okrajů, barev či flexboxu), které se aplikují přímo do HTML/JSX kódu komponent. 

Tento přístup eliminuje nutnost neustálého přepínání mezi oddělenými soubory se styly a aplikační logikou.
Vývojář si tak může při pohledu na zdrojový kód komponenty rovnou představit její rozvržení. 
Zásadní výhodou Tailwindu je přítomnost \enquote{Just-In-Time} kompilátoru, který při sestavování aplikace prohledá kód a do výsledného generovaného CSS souboru zahrne výhradně ty styly, které byly skutečně použity. 
Tím je zajištěna minimální výsledná velikost aplikace. 

\subsection{Shadcn}
Zatímco Tailwind poskytuje nástroje pro základní formátování, pro komplexnější interaktivní prvky uživatelského rozhraní byla využita kolekce komponent Shadcn/ui. 
Nejedná se o tradiční knihovnu, která se instaluje přes \texttt{npm}, ale spíše o sadu přizpůsobitelných komponent generovaných přímo do zdrojového kódu projektu, nad kterými má vývojář plnou kontrolu. 

Komponenty z rodiny Shadcn jsou postavené na knihovně Radix UI, což garantuje jejich vysokou přístupnost (accessibility) pro uživatele s omezením, včetně podpory pro ovládání z klávesnice a čtečky obrazovky. 
V aplikaci byly tyto prvky využity pro tvorbu rozbalovacích nabídek, tlačítek, modálních oken a vyskakovacích upozornění. 

\section{Zpracování a vizualizace formálních modelů}
Hlavní cíl této práce spočívá ve schopnosti aplikace syntakticky analyzovat uživatelem zadané CCS procesy a vizualizovat odpovídající datové struktury. 
K tomuto účelu byly využity specifické knihovny pro zpracování textu a interaktivní vykreslení grafů.

\subsection{Peggy.js}
Pro zpracování zadaných programů v kalkulu komunikujících systémů (CCS) byl využit generátor syntaktických analyzátorů Peggy.js. 
Tento nástroj umožňuje definovat formální gramatiku jazyka pomocí zápisu PEG (Parsing Expression Grammar). 
Z této deklarativní definice Peggy.js následně vygeneruje deterministický a vysoce efektivní parser přímo v jazyce JavaScript.

V aplikaci je tento parser zodpovědný za prvotní kontrolu syntaktické správnosti zadaných CCS výrazů (validace závorek, prefixů, operátorů apod.). 
V případě chybného vstupu dokáže vygenerovat detailní chybová hlášení identifikující řádek a sloupec chyby, čímž uživateli poskytuje okamžitou zpětnou vazbu srovnatelnou s moderními vývojovými prostředími. 

Při úspěšném zpracování transformuje Peggy.js vstupní text do struktury abstraktního syntaktického stromu (AST). 
Tento strom je následně zkontrolován dodatečnými sémantickými kontrolami (např. kontrola existence deklarací všech volaných procesních konstant). 
Po úspěšné validaci tvoří AST základní datovou strukturu pro generování LTS grafu a pro konstruování jednotlivých kroků strukturální operační sémantiky.

\subsection{CodeMirror}
O zadávání a zobrazování zdrojových kódů CCS programů se stará editorová komponenta CodeMirror. 
Jedná se o flexibilní webový textový editor navržený pro implementaci vývojářských nástrojů přímo v prohlížeči. 

Pro potřeby této práce byl vytvořen vlastní režim lexikální analýzy (syntax highlighting) specifický pro jazyk CCS, jehož pravidla odpovídají gramatice použité v Peggy.js. 
Díky tomu jsou vizuálně odděleny akce, operátory, závorky, názvy procesů a speciální tichá akce $\tau$. 
Kromě hlavního okna pro zápis kódu je CodeMirror v aplikaci využit i v režimu pouze pro čtení (read-only) k formátování kratších úseků textu. 
To umožňuje zachovat konzistentní zvýraznění syntaxe v celém uživatelském rozhraní, od instrukcí a nápověd až po samotný text interaktivních SOS důkazů.

\subsection{Cytoscape.js a doplňky}
Pro vizualizaci ohodnocených přechodových systémů a syntaktických stromů, z nichž mohou vzniknout rozsáhlé orientované grafy, byla použita knihovna Cytoscape.js. 
Jedná se o vysoce optimalizovanou JavaScriptovou knihovnu určenou primárně pro modelování a interaktivní vizualizaci komplexních síťových struktur.

V rámci aplikace Cytoscape vykresluje stavy reprezentované jako vrcholy grafu a přechody definované akcemi ve formě orientovaných hran. 
Knihovna poskytuje uživateli možnost graf volně přibližovat, oddalovat a posouvat. Dále umožňuje pokročilou interaktivitu, jako je dynamické propojení syntaktického stromu s textovým editorem CodeMirror, kdy se při najetí myší na vrchol grafu zvýrazní odpovídající část kódu. 

Pro rozšíření možností vizualizace byly využity následující doplňky (pluginy) knihovny:
\begin{itemize}
    \item \textbf{cytoscape-dagre:} Rozšíření, které integruje algoritmus Dagre pro hierarchické a deterministické rozvržení grafu. Tento algoritmus je využit při vizualizaci postupně rozvíjených stavových prostorů LTS.
    \item \textbf{cytoscape-svg (cysvg):} Tento doplněk doplňuje aplikaci o funkcionalitu exportu vykresleného grafu do bezztrátového vektorového formátu SVG.
\end{itemize}

\section{Správa dat}
Jelikož je aplikace z architektonického hlediska navržena bez serveru (serverless klient), neobsahuje tradiční backendovou relační databázi. 
Veškerá data, nastavení i modely jsou ukládány a spravovány výhradně na straně klientského prohlížeče.

\subsection{Formát JSON}
Základním formátem pro strukturu a přenos všech dat v aplikaci je formát JSON (JavaScript Object Notation). 
Tento datový formát je nativně podporován jazykem JavaScript a je jednoduše čitelný jak pro stroj, tak pro člověka. 
V architektuře projektu je použit na několika místech:
\begin{itemize}
    \item Slouží jako základní formát pro ukládání všech vytvořených CCS programů, což umožňuje jejich snadné zálohování či sdílení mezi uživateli.
    \item Využívá se i pro uložení předpřipravených ukázkových vstupů (soubor \texttt{baseProgramList.json}), kde je definována sada validních CCS programů demonstrujících různé příklady z oblasti paralelních systémů.
    \item Poskytuje strukturu pro správu lokalizace v rámci uživatelského rozhraní.
    \item Nabízí dodatečnou možnost strojového exportu vygenerovaných grafů z knihovny Cytoscape, díky čemuž lze s těmito daty dále pracovat v externích analytických nástrojích.
\end{itemize}

\subsection{LocalStorage}
Pro persistenci dat mezi jednotlivými relacemi uživatele (např. po zavření nebo obnovení stránky) je využíváno webové rozhraní \texttt{localStorage}, které je součástí Web Storage API. 
Jedná se o lokální \enquote{key-value} úložiště alokované prohlížečem a izolované čistě pro danou webovou doménu.

Prostřednictvím \texttt{localStorage} aplikace uchovává veškerá uživatelská data, včetně seznamu vytvořených CCS programů, aktuálního nastavení lokalizace nebo preferencí zobrazení. 

\section{Doplňkové nástroje}

\subsection{Lokalizace}
Jelikož se problematika formální verifikace a kalkulu komunikujících systémů může vyučovat v českém i anglickém jazyce, byla aplikace od začátku navržena s podporou pro internacionalizaci (i18n). 
Textové řetězce nejsou pevně zakomponovány v logice komponent, ale jsou externě a dynamicky načítány pomocí lokalizačního frameworku \texttt{i18next}. 
Tato architektura zajišťuje spolehlivý překlad veškerých prvků uživatelského rozhraní, chybových zpráv i nápověd do češtiny a angličtiny, a to se schopností okamžitého přepnutí jazyka za běhu aplikace, bez nutnosti obnovení stránky.

\subsection{Progresivní webová aplikace}
S cílem maximalizovat dostupnost nástroje byla celá aplikace navržena jako progresivní webová aplikace (PWA). 
Tohoto konceptu bylo dosaženo integrací modulu \texttt{vite-plugin-pwa}. 

Architektura PWA využívá Service Workerů, což jsou skripty běžící odděleně na pozadí prohlížeče, které se starají o proaktivní ukládání (caching) statických souborů, knihoven i lokalizačních dat do mezipaměti. 
Aplikace tak po prvním načtení plnohodnotně funguje i bez připojení k internetu, a to včetně kompilace výrazů, tvorby grafů a ukládání programů. 
Součástí konfigurace je i webový manifest, který uživatelům dává možnost \enquote{nainstalovat} si aplikaci přímo do jakéhokoli operačního systému svého zařízení. 
Aplikace se následně spouští ve vlastním vyhrazeném okně bez prvků internetového prohlížeče.

\subsection{ESLint}
Kvalitu, čistotu a jednotnost zdrojového kódu během vývoje pomáhal udržovat nástroj pro statickou analýzu ESLint. 
Tento linter v reálném čase analyzuje napsaný kód a detekuje potenciální syntaktické nedostatky, sémantické chyby či antivzory ještě před samotným sestavením aplikace. 
V kombinaci se striktními konfiguračními pravidly pro TypeScript a knihovnu React tento nástroj průběžně vynucuje osvědčené postupy a omezuje tím vznik běžných programátorských chyb, což v důsledku zvyšuje stabilitu aplikace.

\subsection{Lighthouse}
Při ladění aplikace a přípravě na odevzdání byl pravidelně využíván analytický nástroj Google Lighthouse. 
Jedná se o automatizovaný open-source nástroj určený pro auditování a vylepšování kvality webových stránek. 

Lighthouse simuluje průchod aplikací a generuje detailní metriky pokrývající výkon, SEO a dodržení standardů webových stránek. 
V rámci této práce byl nástroj primárně využit také k identifikaci slabých míst v oblasti přístupnosti (accessibility). 
Na základě vygenerovaných reportů byly v aplikaci provedeny úpravy týkající se barevného kontrastu, sémantiky HTML značek a dostupnosti formulářových prvků pro čtečky obrazovky. 
Cílem bylo maximalizovat skóre v těchto metrikách a zajistit, aby byla aplikace připravená na případné produkční nasazení, kde by již musela splňovat zákon o přístupnosti, již zmíněný v kapitole \ref{sec:funkcni_pozadavky}.